% ============================================================
% Electronic Companion for:
% "When Front-Stage Innovation Overloads Back-Stage Workers"
% Submitted to M&SOM
% ============================================================

\documentclass[msom,dblanonrev]{informs4}
\usepackage{eqndefns-left}
\RequirePackage{tgtermes}
\RequirePackage{newtxtext}
\RequirePackage{newtxmath}
\RequirePackage{bm}
\RequirePackage{endnotes}

\DoubleSpacedXI

% Natbib setup for author-date style
\usepackage{natbib}
\bibpunct[, ]{(}{)}{,}{a}{}{,}%
\def\bibfont{\small}%
\def\bibsep{\smallskipamount}%
\def\bibhang{24pt}%
\def\newblock{\ }%
\def\BIBand{and}%

% Setup of equation and theorem numbering
\EquationsNumberedThrough
\TheoremsNumberedThrough
\ECRepeatTheorems

% Additional packages
\usepackage{booktabs}
\usepackage{multirow}

% Notation
\newcommand{\Gp}{[G]^{+}}
\newcommand{\Gm}{[-G]^{+}}
\newcommand{\Keff}{K_{\mathrm{eff}}}

% Manuscript number
\MANUSCRIPTNO{}

\begin{document}

% Running author and title
\RUNAUTHOR{Authors suppressed}
\RUNTITLE{Front-Stage Innovation Overloads Back-Stage Workers -- EC}

% Full title
\TITLE{Electronic Companion to ``When Front-Stage Innovation Overloads Back-Stage Workers''}

% Authors and affiliations
\ARTICLEAUTHORS{%
\AUTHOR{Authors suppressed for double-blind review}
\AFF{Affiliations suppressed, \EMAIL{suppressed}}
}

% Abstract
\ABSTRACT{This electronic companion provides full proofs, robustness analyses, and calibration details for the main paper.}%

\KEYWORDS{electronic companion}

\maketitle

\section{Full Proofs}\label{sec:ec:proofs}

We restate each result for convenience and provide complete proofs. The model equations are:
\begin{align}
D &= (1 + aT) \cdot F^b, \label{ec:demand} \\
G &= D - K \cdot h(W), \quad h(W) = (1-p) + pW, \label{ec:gap} \\
\frac{dK}{dt} &= r\,\Gp\, W - \delta K + u, \label{ec:dK} \\
\frac{dW}{dt} &= (\Gm + f)(1-W) - (\mu + \sigma\,\Gp)\,W, \label{ec:dW} \\
\frac{dT}{dt} &= \beta\,\Gm\,(1-T) - \eta\beta\,\Gp\, T - \lambda T. \label{ec:dT}
\end{align}
Customer lifetime value is denoted CLV throughout; starred variables ($K^*$, $W^*$, $T^*$) denote steady-state equilibrium values.

\subsection{State-Space Invariance}\label{sec:ec:invariance}

\begin{lemma}[Forward Invariance]\label{lem:invariance}
The set $\Omega = \{(K, W, T) : K > 0,\; W \in [0,1],\; T \in [0,1]\}$ is forward-invariant under the flow of the ordinary differential equation (ODE) system \eqref{ec:dK}--\eqref{ec:dT} for all $t \geq 0$, provided $u > 0$, $f \geq 0$, and all parameters are non-negative.
\end{lemma}

\begin{proof}
At each boundary, the vector field points into $\Omega$: $\dot W|_{W=0} = \Gm + f \geq 0$; $\dot W|_{W=1} = -(\mu + \sigma\Gp) \leq 0$; $\dot T|_{T=0} = \beta\Gm \geq 0$; $\dot T|_{T=1} = -(\eta\beta\Gp + \lambda) \leq 0$; $\dot K|_{K=0} = r\Gp W + u > 0$. Forward invariance follows by the standard barrier argument.\Halmos
\end{proof}

\begin{remark}[Piecewise-Smooth Dynamics]\label{rem:pwsmooth}
The system \eqref{ec:dK}--\eqref{ec:dT} is piecewise-smooth, with a non-smooth gradient at the switching manifold $\mathcal{M} = \{G = 0\}$. Existence and uniqueness of solutions follow from Filippov's (\citeyear{filippov1988differential}) framework, provided trajectories cross $\mathcal{M}$ transversally; numerical verification at both equilibria confirms this under baseline parameters. The implicit function theorem arguments underpinning the analytical results are applied at surplus equilibria where $G^* = -S^* < 0$, so differentiability is guaranteed by the smooth structure of the ODE within $\{G < 0\}$; smoothing robustness is verified in \S\ref{sec:ec:smoothing}.
\end{remark}

\begin{remark}[Computability of $F_{\mathrm{crit}}$]\label{rem:Fcrit_implicit}
The implicit balance condition $K^*h(W^*)=(1+aT^*)F^b+S^*$ rearranges to $F=(K^*h(W^*)/(1+aT^*))^{1/b}$, but this is self-referential since $W^*$ and $T^*$ depend on $F$. The computable form evaluates $h$ and the demand term at the boundary $S^*\to 0$ where $T^*\to 0$, yielding the closed-form threshold $F_{\mathrm{crit}} = (K^* h(f/(f+\mu)))^{1/b}$.
\end{remark}

\subsection{Equilibrium Existence and Uniqueness}\label{sec:ec:existence}

Before proving the main theorems, we establish that the surplus equilibrium invoked by Theorems~1--3 exists and is unique. The equilibrium expressions derived in \S4.1 of the main paper are conditional on a surplus $S^* > 0$, but $S^*$ itself depends on the equilibrium values $W^*$ and $T^*$ through $S^* = K^* h(W^*) - D(T^*, F)$, creating a fixed-point problem. We resolve this by reducing the system to a single implicit equation in $S^*$ and establishing that it has exactly one positive solution.

\begin{lemma}[Surplus Equilibrium Existence; Uniqueness under C2]\label{lem:existence}
For any $F < F_{\mathrm{crit}}$, a surplus equilibrium $(K^*, W^*, T^*)$ with $S^* > 0$ exists. Under Regularity Condition C2 (single-crossing condition C2-SC: $\beta(f+\mu)<\lambda$; analytic, see Step~1 of the Theorem~1 proof, EC~\S EC.1.3), this equilibrium is unique.
\end{lemma}

\begin{proof}
At surplus equilibrium, $K^* = u/\delta$. The equilibrium conditions yield $W^*(S) = (S{+}f)/(S{+}f{+}\mu)$ and $T^*(S) = \beta S/(\beta S{+}\lambda)$, both continuous and strictly increasing in surplus $S$. Define the surplus residual $\Phi(S) = (u/\delta)[(1{-}p) + pW^*(S)] - (1{+}aT^*(S))F^b - S$. Then $\Phi(0) > 0$ when $F < F_{\mathrm{crit}}$, and $\Phi(S) \to -\infty$ as $S\to\infty$. Existence follows by IVT. For uniqueness, $\Phi'(S) = \tau_1(S) - \tau_2(S) - 1$ where $\tau_1, \tau_2 > 0$ are decreasing. Under C2-SC ($\beta(f{+}\mu) < \lambda$), the equation $\Phi'(S) = 0$ has at most one solution (a decreasing function equals an increasing function), so $\Phi$ has exactly one positive root. C2-SC holds for all $f \in [0, 0.10]$: $\beta(f{+}\mu) \leq 0.018 < 0.025 = \lambda$.\Halmos
\end{proof}

\begin{remark}[Low-Trust Interior Attractor]\label{rem:low_trust}
A second fixed point exists for large $F$, with $T^*{=}0$ and $W^*_{\mathrm{low}} = f/(f{+}\mu{+}\sigma G^*)$. The Jacobian has strictly negative eigenvalues (verified numerically for $F \in [1.0, 1.5]$). This interior attractor, jointly with the surplus equilibrium of Lemma~\ref{lem:existence}, establishes the bistability invoked by Theorem~3.
\end{remark}

\begin{lemma}[Forward-Boundedness of Trajectories]\label{lem:bounded}
For any finite $F$ and any initial condition $(K_0, W_0, T_0) \in \mathbb{R}_{>0} \times [0,1] \times [0,1]$, the solution satisfies $K(t) \leq \max(K_0,\, C(F)/\delta)$ for all $t \geq 0$, where $C(F) = r(1+a)F^b + u$. Combined with $W(t)\in[0,1]$ and $T(t)\in[0,1]$ (both preserved by the ODE structure), the state space is forward-invariant and bounded.
\end{lemma}

\begin{proof}
Under overload ($G > 0$), we bound $G^+ = D(T,F) - K\,h(W) \leq D(T,F) \leq (1+a)F^b$ (since $T \leq 1$) and $W \leq 1$. Therefore
\begin{align*}
  \dot{K} &= r G^+ W - \delta K + u \\
    &\leq r(1+a)F^b - \delta K + u = C(F) - \delta K.
\end{align*}
By Gronwall's lemma, $K(t) \leq \max(K_0,\, C(F)/\delta)$ for all $t \geq 0$. Under surplus ($G \leq 0$), $G^+=0$ and $\dot{K} = -\delta K + u \leq 0$ whenever $K \geq u/\delta$, so $K$ is bounded above by $\max(K_0, u/\delta)$. In both regimes the upper bound $\max(K_0, C(F)/\delta)$ holds. At baseline parameters with $F = 1.2$, this gives $K_{\max} \approx 25.3$, far above the operating range $K^* = 1.67$. Since $W$ and $T$ are bounded to $[0,1]$ by construction, the full state space is forward-invariant and compact for all finite $F$. \Halmos
\end{proof}

\begin{remark}[Scope of global convergence]\label{rem:global_conv}
Lemma~\ref{lem:bounded} ensures all trajectories remain in a bounded compact set. Which attractor a trajectory reaches depends on its initial condition relative to the basin boundary (Figure~2, main paper). For $F \geq F_{\mathrm{crit}}$, only the low-trust attractor is stable (Remark~\ref{rem:low_trust}). A complete global Lyapunov proof is left to future work; the principal obstacle is the switching manifold $\{G{=}0\}$ where the standard LaSalle argument does not apply directly, requiring Filippov-specific Lyapunov construction~\citep{filippov1988differential}.
\end{remark}

\begin{remark}[Maintained Conditions]\label{rem:conditions}
Four conditions are maintained throughout: (C1)~$\Phi(0) > 0$ (i.e., $F < F_{\mathrm{crit}}$; analytic via IVT); (C2)~uniqueness of $S^*$, established analytically under C2-SC ($\beta(f+\mu)<\lambda$, which holds for all $f\in[0,0.10]$ independent of $b$, $p$, $\sigma$; see Step~1 of the Theorem~1 proof, EC~\S EC.1.3); (C3)~regularity $\Phi'(S^*) < 0$, a consequence of~(C1)--(C2), required for the implicit function theorem in Theorems~1--2; and (C4)~low-trust attractor existence $F^b > (u/\delta)\,h(W^*_{\mathrm{low}})$, verified numerically (Jacobian eigenvalues across $F \in [1.0, 1.5]$).
\end{remark}

\begin{remark}[Bifurcation Classification at $F_{\mathrm{crit}}$]\label{rem:bifurcation}
The transition at $F_{\mathrm{crit}}$ is a saddle-node (fold) bifurcation in the surplus variable $S$: the stable surplus equilibrium and the separating saddle coalesce as $F \nearrow F_{\mathrm{crit}}$. The $(W,T)$ sub-system eigenvalues retain strictly negative real parts throughout $F \in [0, F_{\mathrm{crit}}]$ (baseline: complex-conjugate pair with $\mathrm{Re} \approx -0.07$; characteristic polynomial constant term $>0$ up to $F_{\mathrm{crit}}$), ruling out Hopf bifurcation.
\end{remark}

\subsection{Proof of Theorem 1 (Front-Stage Trap)}\label{sec:ec:proof1}

\begin{theorem}[Front-Stage Trap, restated]
Suppose $b > 1$ and $p > 0$. Define the threshold $F_{\mathrm{crit}} = \bigl((u/\delta)\,h(f/(f+\mu))\bigr)^{1/b}$. Then: (i)~for $F < F_{\mathrm{crit}}$, the unique surplus equilibrium exists with $T^* > 0$ and $\mathrm{CLV}^* > \mathrm{CLV}_{\mathrm{low}}$; (ii)~for $F \geq F_{\mathrm{crit}}$, no surplus equilibrium exists; all forward trajectories remain bounded (Lemma~\ref{lem:bounded}) and converge numerically to the low-trust attractor with $\mathrm{CLV} \approx \mathrm{CLV}_{\mathrm{low}}$.
\end{theorem}

\begin{proof}
We show that the surplus equilibrium exists if and only if $F < F_{\mathrm{crit}}$, and that beyond this threshold the system is attracted to the low-trust fixed point.

\textbf{Step 1: Surplus equilibrium requires $\Phi(0) > 0$.} From Lemma~\ref{lem:existence}, the surplus equilibrium exists if and only if $\Phi(0) = (u/\delta)\,h(f/(f+\mu)) - F^b > 0$, which rearranges to $F < F_{\mathrm{crit}} = \bigl((u/\delta)\,h(f/(f+\mu))\bigr)^{1/b}$. (Here we use the fact that at $S^* = 0$, $T^* \to 0$ and $W^* \to f/(f+\mu)$.)

\textbf{Step 2: Within the surplus regime, surplus decreases in $F$.} At the surplus equilibrium, $K^* = u/\delta$ is independent of $F$. Differentiating the surplus residual, $dS^*/dF < 0$ for $F$ near $F_{\mathrm{crit}}$ because the demand derivative $b(1+aT^*)F^{b-1}$ dominates the indirect effects through $W^*$ and $T^*$. As $F \to F_{\mathrm{crit}}$, $S^* \to 0$.

\textbf{Step 3: Beyond $F_{\mathrm{crit}}$, overload drives collapse.} For $F \geq F_{\mathrm{crit}}$, $G > 0$, and from (\ref{ec:dT}):
\[
\frac{dT}{dt} = -T(\eta\beta G + \lambda) < 0 \quad \text{for all } T > 0.
\]
Trust declines monotonically. Through Loop R2 ($W\!\downarrow \to h(W)\!\downarrow \to G\!\uparrow$), the gap widens, accelerating trust decline. The convexity of demand in $F$ ($\partial^2 D/\partial F^2 > 0$ for $b > 1$) ensures the gap widens super-linearly, making the collapse self-reinforcing under Loop R2; convergence to the low-trust attractor is confirmed numerically across the studied parameter configurations (Tier~C; main paper Remark~1).% NB: hardcoded remark number; update if main-paper Remark (Analytical Transparency) is renumbered. The system converges to the low-trust interior attractor (Remark~\ref{rem:low_trust}).

\textbf{Boundary conditions.} When $b \leq 1$: demand is concave in $F$; the balancing loop B1 can compensate and the regime shift does not occur. When $p = 0$: $h(W) = 1$, severing Loop R2; the system may still experience trust decline through the direct overload channel but without the workforce-mediated amplification.
\end{proof}


\subsection{Proof of Theorem 2 (Good Jobs Buffer)}\label{sec:ec:proof2}

\begin{theorem}[Good Jobs Buffer, restated]
Suppose $p > 0$ and the surplus equilibrium is unique and stable (Regularity Condition: $\Phi'(S^*) < 0$; Lemma~\ref{lem:existence}). Then $\partial\,\mathrm{CLV}^*/\partial f > 0$ for all $f \geq 0$.
\end{theorem}

\begin{proof}
We show that each link in the chain $f \to W^* \to h(W^*) \to \Keff \to S^* \to T^* \to \mathrm{CLV}^*$ is strictly positive.

\textbf{Link 1: $\partial W^*/\partial f > 0$.} At the surplus equilibrium, from (\ref{ec:dW}) with $\Gp = 0$:
\begin{align*}
0 &= S^*(1-W^*) + f(1 - W^*) - \mu W^* \\
  &= (S^*+f)(1-W^*) - \mu W^*.
\end{align*}
Solving: $W^* = (S^* + f)/(S^* + f + \mu)$.
Differentiating: $\partial W^*/\partial f = \mu / (S^* + f + \mu)^2 > 0$. This is positive, decreasing in $f$ (diminishing returns), and increasing in $\mu$ (higher natural decay amplifies the marginal benefit of investment).

\textbf{Link 2: $\partial h / \partial W = p > 0$.} By construction of $h(W) = (1-p) + pW$.

\textbf{Link 3: $\partial S^*/\partial h = K^* > 0$.} Since $S^* = K^* h(W^*) - D(T^*, F)$ and $K^* = u/\delta > 0$.

\textbf{Link 4: $\partial T^*/\partial S^* > 0$.} From the equilibrium expression $T^* = \beta S^* / (\beta S^* + \lambda)$: $\partial T^*/\partial S^* = \beta\lambda / (\beta S^* + \lambda)^2 > 0$.

\textbf{Link 5: $\partial\,\mathrm{CLV}/\partial T^* > 0$.} By assumption on the CLV functional form ($w$ and $\rho$ both increasing in $T$).

\textbf{Combining links via the chain rule:}
\[
\frac{\partial\,\mathrm{CLV}^*}{\partial f} = \underbrace{\frac{\partial\,\mathrm{CLV}}{\partial T^*}}_{>0} \cdot \underbrace{\frac{\partial T^*}{\partial S^*}}_{>0} \cdot \underbrace{K^*}_{>0} \cdot \underbrace{p}_{>0} \cdot \underbrace{\frac{\partial W^*}{\partial f}}_{>0} > 0.
\]
Diminishing returns: $\partial^2\,\mathrm{CLV}^*/\partial f^2 < 0$ follows from the concavity of $W^*$ in $f$ (the $(S^*+f+\mu)^{-2}$ factor is decreasing in $f$) propagated through the chain.

\textbf{General-equilibrium correction.} The full equilibrium embeds a feedback loop $f \to W^* \to h \to S^* \to T^* \to D$, which partially offsets the direct surplus gain by raising demand through higher trust. To establish that this feedback does not reverse the sign, we apply the implicit function theorem directly.

\textit{Sign established analytically.} The equilibrium condition $\Phi(S^*) = 0$ defines $S^*$ implicitly as a function of $f$. Differentiating:
\[
\frac{dS^*}{df} = -\frac{\partial\Phi/\partial f}{\Phi'(S^*)}.
\]
Since $\partial\Phi/\partial f = (u/\delta)\,p\,\mu/(S^*\!+\!f\!+\!\mu)^2 > 0$ and $\Phi'(S^*) < 0$ (Regularity Condition C3), we have $dS^*/df > 0$ for all $f \geq 0$. The sign of $d\,\mathrm{CLV}^*/df$ is therefore analytically established; it does not depend on the magnitude of the trust-demand feedback.

\textit{Magnitude of the GE correction (numerically confirmed).} The trust-demand feedback damps the surplus gain by a factor proportional to $aF^b\,\beta\lambda/(\beta S^*+\lambda)^2 / |\Phi'(S^*)|$, which we refer to as the GE correction ratio. At baseline parameters this ratio is approximately 0.09, confirmed across 100 parameter configurations spanning the empirically relevant range. The correction dampens but does not reverse the direct effect; $d\,\mathrm{CLV}^*/df > 0$ throughout.
\end{proof}


\subsection{Proof of Theorem 3 (Burnout Cliff)}\label{sec:ec:proof3}

\begin{theorem}[Burnout Cliff, restated]
Suppose $\sigma > 1$, $p > 0$, and $F_0 < F_{\mathrm{crit}}$ (a surplus equilibrium exists at $F_0$). Consider a firm at the surplus equilibrium that implements a step increase in $F$ from $F_0$ to $F_0 + \Delta F$. The outcome is governed by two thresholds:
\begin{enumerate}
    \item[(i)] \emph{Safety threshold (permanent-collapse onset).} For
    \begin{align*}
       \Delta F &\;\geq\; \Delta F_{\mathrm{collapse}} = F_{\mathrm{crit}}(f) - F_0 \\
         &= \Bigl(\frac{u}{\delta}\,h\!\Bigl(\frac{f}{f+\mu}\Bigr)\Bigr)^{\!1/b} \!- F_0,
    \end{align*}
    no surplus equilibrium exists at the post-upgrade level; the low-trust attractor is the only locally stable state (Remark~\ref{rem:low_trust}), and convergence to it is confirmed numerically across the studied initial conditions. For $\Delta F < \Delta F_{\mathrm{collapse}}$, a new surplus equilibrium exists at the post-upgrade level. Convergence to this equilibrium is basin-dependent; numerical verification confirms recovery from pre-shock equilibrium initial conditions across the parameter space studied (EC~\S EC.4.2).
    \item[(ii)] \emph{Speed diagnostic (instantaneous-overload onset).} A secondary threshold
    \[
       \Delta F_{\mathrm{rapid}} \;=\; \frac{S_0}{b\,F_0^{b-1}(1+aT^*)}
    \]
    marks the boundary above which the service gap turns positive at $t = 0^+$, initiating a rapid self-reinforcing collapse to the low-trust attractor. Because $\Delta F_{\mathrm{rapid}} > \Delta F_{\mathrm{collapse}}$ at baseline ($0.40$ vs.\ $0.31$), a firm upgrading by $\Delta F \in [\Delta F_{\mathrm{collapse}}, \Delta F_{\mathrm{rapid}})$ collapses permanently but through a \emph{slow} trajectory invisible on a standard fiscal-quarter horizon. The safety threshold is $\Delta F_{\mathrm{collapse}}$; $\Delta F_{\mathrm{rapid}}$ characterises collapse speed, not safety.
\end{enumerate}
\end{theorem}

\begin{proof}
\textbf{Step 1: Two thresholds.} The permanent-collapse threshold follows directly from Theorem~1: for $F_0 + \Delta F \geq F_{\mathrm{crit}}$, the surplus residual $\Phi(S)$ (Lemma~\ref{lem:existence}) at $S=0$ gives $\Phi(0) \leq 0$, so no surplus equilibrium exists. Therefore
\[
   \Delta F_{\mathrm{collapse}} = F_{\mathrm{crit}} - F_0
     = \Bigl(\frac{u}{\delta}\,h\!\Bigl(\frac{f}{f+\mu}\Bigr)\Bigr)^{1/b} - F_0.
\]
At baseline ($F_0=0.80$): $\Delta F_{\mathrm{collapse}} = 1.1104 - 0.80 = 0.310$.

For $\Delta F < \Delta F_{\mathrm{collapse}}$, the post-upgrade $F_0+\Delta F < F_{\mathrm{crit}}$, so a new surplus equilibrium exists (Lemma~\ref{lem:existence}). Convergence to this equilibrium is basin-dependent; Remark~\ref{rem:global_conv} establishes boundedness, and numerical verification confirms convergence from pre-shock equilibrium initial conditions across the parameter configurations studied.

\textbf{Step 2: Instantaneous-overload diagnostic.} The first-order demand increment after the upgrade is $\Delta D \approx (1+aT^*)\,b\,F_0^{b-1}\,\Delta F$. Setting $\Delta D = S_0$ (the pre-upgrade surplus):
\begin{equation}\label{ec:DFrapid}
   \Delta F_{\mathrm{rapid}} = \frac{S_0}{b\,F_0^{b-1}(1+aT^*)}.
\end{equation}
At baseline: $\Delta F_{\mathrm{rapid}} = 0.401$. Since $\Delta F_{\mathrm{rapid}} > \Delta F_{\mathrm{collapse}}$ (0.401 vs.\ 0.310), the interval $[\Delta F_{\mathrm{collapse}},\Delta F_{\mathrm{rapid}})$ defines a \emph{slow-collapse zone}: the surplus equilibrium is gone, but $G(t_0^+) < 0$, so collapse proceeds through a gradual erosion of surplus rather than immediate overload. For $\Delta F \geq \Delta F_{\mathrm{rapid}}$, $G(t_0^+) > 0$ and collapse is rapid.

\textbf{Step 3: Dynamics and timescales.} Under overload ($G > 0$), the tipping threshold $W_{\mathrm{tip}} = f/(f+\mu+\sigma G_0^+)$ marks the level below which depletion exceeds recovery. Two timescales govern speed: $\tau_{\mathrm{adapt}} = 1/r \approx 2.5$ mo (Loop B1) and $\tau_{\mathrm{deplete}} = 1/(\mu+\sigma G_0^+) \in [1.5,4]$ mo (Loop R3). Near $\Delta F_{\mathrm{collapse}}$, convergence to the low-trust attractor may require hundreds of months.

\textbf{Step 4: Three zones.} $\Delta F_{\mathrm{collapse}}$ and $\Delta F_{\mathrm{rapid}}$ define: (i)~\emph{safe zone} ($\Delta F < \Delta F_{\mathrm{collapse}}$), where a surplus equilibrium exists; (ii)~\emph{slow-collapse zone} ($\Delta F_{\mathrm{collapse}} \leq \Delta F < \Delta F_{\mathrm{rapid}}$), where $G(t_0^+) < 0$ but no surplus equilibrium exists, so collapse is gradual; (iii)~\emph{rapid-collapse zone} ($\Delta F \geq \Delta F_{\mathrm{rapid}}$), where $G(t_0^+) > 0$ immediately. The ordering $\Delta F_{\mathrm{rapid}} > \Delta F_{\mathrm{collapse}}$ reverses at high $f$ ($f{=}0.10$: $\Delta F_{\mathrm{collapse}} = 0.49 > \Delta F_{\mathrm{rapid}} \approx 0.40$). Simulations ($t_{\mathrm{end}}{=}500$, collapse criterion $T(\infty) < 0.15$) confirm collapse boundaries across three workforce scenarios at $F_0{=}0.80$. Minor discrepancies between simulated and analytic thresholds reflect slow-convergence transients near the boundary and basin-boundary effects. The correct safety threshold is $\Delta F_{\mathrm{collapse}}$; $\Delta F_{\mathrm{rapid}}$ characterises collapse speed, not safety. \Halmos
\end{proof}


\subsection{Proof of Proposition 4 (Complementarity)}\label{sec:ec:proof4}

\setcounter{proposition}{3}
\begin{proposition}[Digitalisation-Workforce Complementarity, restated]
$\partial \Keff / \partial W = pK > 0$ and $\partial^2 \Keff / \partial W \partial K = p > 0$. Workforce well-being determines $F_{\max}(W) = (K^* h(W))^{1/b}$, which is strictly increasing in $W$.
\end{proposition}

\begin{proof}
\textbf{Part 1: Cross-partial.} From $\Keff = K \cdot [(1-p) + pW]$: $\partial \Keff / \partial W = pK$ and $\partial^2 \Keff / \partial W \partial K = p$. Both are strictly positive for $p > 0$.

\textbf{Part 2: $F_{\max}(W)$ at the regime boundary.} $F_{\max}(W)$ is defined as the value of $F$ at which the surplus equilibrium of Lemma~\ref{lem:existence} ceases to exist, i.e., at which $S^* = 0$. From Lemma~\ref{lem:existence}, at $S^* = 0$:
\[
T^*(S) = \frac{\beta S}{\beta S + \lambda} \;\xrightarrow{S \to 0}\; 0.
\]
Therefore, substituting $T^* = 0$ into the $F_{\mathrm{crit}}$ formula from Theorem~1 (Eq.~\ref{ec:demand} at $S^* = 0$):
\begin{align*}
F_{\max}(W) &= \Bigl(K^* \cdot [(1-p) + pW]\Bigr)^{1/b} \\
  &= (K^*\,h(W))^{1/b}.
\end{align*}
This boundary formula is exact at $S^* = 0$ because $T^* \to 0$ is a direct consequence of the equilibrium structure, not an approximation.

\textbf{Full implicit differentiation at interior equilibria ($S^* > 0$).} At interior equilibria, $T^* > 0$ and $F_{\max}$ solves the fixed-point equation $(K^* h(W))/(1 + aT^*(W,F_{\max})) = F_{\max}^b$. Total differentiation with respect to $W$ yields:
\[
\frac{dF_{\max}}{dW} = \frac{1}{b}\,F_{\max}^{1-b}\,\frac{d}{dW}\!\left[\frac{K^* h(W)}{1 + aT^*}\right],
\]
where
\[
\frac{d}{dW}\!\left[\frac{K^* h(W)}{1 + aT^*}\right] = \frac{K^* p}{1+aT^*} - \frac{K^* h(W)\,a}{(1+aT^*)^2}\,\frac{dT^*}{dW}.
\]
Since $T^*(S^*)$ is increasing in $S^*$ and $dS^*/dW > 0$ (Theorem~2, Link~2--Link~3), the correction term $-K^* h(W) a\,(dT^*/dW)/(1+aT^*)^2$ is negative, meaning the interior-equilibrium slope $dF_{\max}/dW$ is somewhat smaller than the boundary-formula slope $(p/b)\,(K^*)^{1/b}[h(W)]^{1/b-1}$. However, the sign remains positive: $dF_{\max}/dW > 0$ for all $p > 0$, $K^* > 0$. The boundary formula therefore provides a conservative upper bound on the interior slope: the interior slope is smaller (by approximately 8\% at baseline parameters), so the formula overstates how rapidly $F_{\max}$ increases with $W$ at positive trust levels. \Halmos
\end{proof}


\subsection{Proof of Corollary 5 (Budget-Constrained Trap)}\label{sec:ec:proof5}

\begin{proof}
Let $B_{\min} \equiv c_F \cdot F_{\mathrm{crit}}$ denote the budget at which the all-$F$ corner exactly attains the Trap threshold.

\textit{Case 1: $B \geq B_{\min}$.} The all-$F$ corner allocates $F = B/c_F \geq F_{\mathrm{crit}}$, where Theorem~1 implies $\partial\,\mathrm{CLV}/\partial F \leq 0$ (the surplus equilibrium has ceased to exist and the system is at or beyond the Trap threshold). Theorem~2 simultaneously guarantees $\partial\,\mathrm{CLV}^*/\partial f > 0$ for all $f \geq 0$. Therefore, reallocating the marginal dollar from $F$ to $f$ produces a strict first-order gain in CLV. The all-$F$ corner is strictly dominated and the optimum satisfies $f^* > 0$ and $F^* < B/c_F$.

\textit{Case 2: $B < B_{\min}$.} The all-$F$ corner allocates $F = B/c_F < F_{\mathrm{crit}}$, so the Trap is not binding and a surplus equilibrium exists. Theorem~2 guarantees $\partial\,\mathrm{CLV}^*/\partial f > 0$, so the all-$F$ corner ($f = 0$) forgoes a positive return to workforce investment.

\textit{Near-$B_{\min}$ sub-case (analytic).} The per-dollar returns to $F$ and $f$ are continuous in $F$ on $\{F < F_{\mathrm{crit}}\}$: both follow from the explicit equilibrium formulae for $W^*(S)$, $T^*(S)$, and Theorems~1--2 via the implicit function theorem. At $B = B_{\min}$ (i.e., $F = F_{\mathrm{crit}}$), Case~1 establishes $\partial\,\mathrm{CLV}/\partial F = 0$ analytically while $\partial\,\mathrm{CLV}^*/\partial f > 0$ (Theorem~2). By continuity, strict dominance of $f$ over $F$ per dollar persists throughout a neighbourhood $B \in (B_{\min}-\varepsilon, B_{\min})$ for some $\varepsilon > 0$.

\textit{Interior sub-case (computational).} For $B \in [B_{\min}/2,\, B_{\min}-\varepsilon)$, strict dominance of $f$ holds computationally across cost ratios $c_F/c_f \in [0.5, 5]$ (the range spanning realistic investment cost structures).

In Case~1 and near $B_{\min}$ in Case~2, $f^* > 0$ follows analytically; the interior of Case~2 is a computational finding. \Halmos
\end{proof}


\subsection{Smoothing Robustness}\label{sec:ec:smoothing}

Table~\ref{tab:ec:smooth} verifies that replacing the exact piecewise-smooth operators $[\cdot]^+$ with a tanh-smoothed approximation shifts $F_{\mathrm{crit}}$ by less than 1\% for $\varepsilon \leq 0.01$, confirming that all three theorems are invariant to smoothing.

\begin{table}
\TABLE
{Sensitivity of $F_{\mathrm{crit}}$ to tanh smoothing bandwidth $\varepsilon$.\label{tab:ec:smooth}}
{\begin{tabular}{@{}lcc@{}}
\toprule
Smoothing $\varepsilon$ & $F_{\mathrm{crit}}$ & Change vs.\ exact (\%) \\
\midrule
Exact ($\varepsilon = 0$) & 1.110 & -- \\
$\varepsilon = 0.001$ & 1.106 & $-0.09$ \\
$\varepsilon = 0.005$ & 1.104 & $-0.27$ \\
$\varepsilon = 0.010$ & 1.100 & $-0.63$ \\
$\varepsilon = 0.050$ & 1.078 & $-2.62$ \\
\bottomrule
\end{tabular}}
{}
\end{table}



\section{Structural Mapping: WCT and Oliva--Sterman (2001)}\label{sec:ec:nesting}

The WCT model recovers the core dynamics of \citeauthor{oliva2001cutting} (\citeyear{oliva2001cutting}) as a special case under three parameter restrictions: (i) $p = 0$ (workforce does not affect $\Keff$, so $h(W)=1$); (ii) $a = 0$ (no trust-demand feedback, so $D = F^b$); and (iii) $F$ constant (no cross-domain dynamics). Under these settings, the trust ODE remains active but effectively decouples from demand; setting additionally $\beta \to \infty$ and $\lambda = 0$ collapses $T$ into the demand function, recovering the O-S perceived-quality attractiveness mechanism. Table~\ref{tab:ec:mapping} provides the term-by-term correspondence.

\begin{table}
\TABLE
{Structural mapping between WCT (restricted) and Oliva--Sterman (2001).\label{tab:ec:mapping}}
{\begin{tabular}{@{}lll@{}}
\toprule
\textbf{WCT term} & \textbf{O\&S counterpart} & \textbf{Correspondence} \\
\midrule
$r\,\Gp\,W$ & Learning from failures & Product of gap and workforce effort \\
$\delta K$ & Depreciation & Constant-rate technology decay \\
$(\Gm + f)(1-W)$ & Recovery + staffing & Surplus and investment raise effort; saturates at $W{=}1$ \\
$(\mu + \sigma\,\Gp)\,W$ & Corner-cutting + fatigue & Overload and decay deplete effort; vanishes at $W{=}0$ \\
\addlinespace
Trust stock $T$ & No O\&S counterpart & Collapses to demand at $\beta \to \infty$, $\lambda = 0$ \\
\bottomrule
\end{tabular}}
{}
\end{table}

The $(1-W)$ saturation factor replaces O-S's nonlinear lookup table; qualitative behaviour is identical. With the three restrictions lifted, the model generates Loops R1, R2, and the cross-domain Front-Stage Trap.

\section{Sensitivity Analysis}\label{sec:ec:sensitivity}

Table~\ref{tab:ec:sensitivity} reports how each theorem's qualitative result responds to $\pm 50\%$ perturbations of each parameter from its baseline value.

\begin{table}
\TABLE
{Sensitivity of theorem predictions to $\pm 50\%$ parameter perturbations. $^{\dagger}$Slower sensing ($r\downarrow$) allows overload more time to erode $W$ before the self-correction loop activates, steepening the cliff.\label{tab:ec:sensitivity}}
{\begin{tabular}{@{}lccc@{}}
\toprule
Parameter perturbation & Thm 1 (Trap) & Thm 2 (Buffer) & Thm 3 (Cliff) \\
\midrule
$b$: 1.3 $\to$ 1.0 & Disappears & Holds & Weakens \\
$b$: 1.3 $\to$ 1.6 & Strengthens & Holds & Strengthens \\
$a$: 0.4 $\to$ 0.2 & Weakens & Holds & Holds \\
$a$: 0.4 $\to$ 0.6 & Strengthens & Holds & Holds \\
$p$: 0.5 $\to$ 0.25 & Weakens & Weakens & Weakens \\
$p$: 0.5 $\to$ 0.75 & Strengthens & Strengthens & Strengthens \\
$\sigma$: 2.0 $\to$ 1.0 & Weakens & Holds & Weakens \\
$\sigma$: 2.0 $\to$ 3.0 & Strengthens & Holds & Strengthens \\
$\eta$: 2.0 $\to$ 1.0 & Weakens & Holds & Holds \\
$\eta$: 2.0 $\to$ 3.0 & Strengthens & Holds & Holds \\
$r$: 0.4 $\to$ 0.2 & Holds & Holds & Strengthens$^{\dagger}$ \\
$r$: 0.4 $\to$ 0.8 & Holds & Holds & Weakens \\
$\delta$: 0.03 $\to$ 0.015 & Weakens & Strengthens & Weakens \\
$\delta$: 0.03 $\to$ 0.06 & Strengthens & Weakens & Strengthens \\
\bottomrule
\end{tabular}}
{}
\end{table}

\textbf{Key observations.} $p$ is the most uniformly influential parameter: reducing it weakens all three theorems simultaneously. $b$ is critical for Theorem~1 (at $b = 1$ the Trap disappears); $\sigma$ for Theorem~3 (at $\sigma = 1$ bistability vanishes). No single perturbation kills all three results.

\begin{table}
\TABLE
{Supplementary sensitivity: leading-indicator lag and direct-channel share across key parameter ranges.\label{tab:ec:sensitivity2}}
{\begin{tabular}{@{}llccc@{}}
\toprule
\textbf{Outcome} & \textbf{Parameter varied} & \textbf{Low} & \textbf{Baseline} & \textbf{High} \\
\midrule
\multirow{2}{*}{Leading-indicator lag (months)$^{\dagger}$}
  & $\beta \in \{0.06, 0.12, 0.24\}$, $\sigma = 2.0$ & 21 & 12 & 6 \\
  & $\sigma \in \{1.0, 2.0, 3.0\}$, $\beta = 0.12$   & 11 & 12 & 12 \\[4pt]
\multirow{1}{*}{Indirect-channel share (\%)}
  & $p \in \{0.25, 0.50, 0.75\}$     & $<$1 & $<$1 & $<$1 \\
\bottomrule
\end{tabular}}
{$^{\dagger}$Leading-indicator lag: months by which $W$ leads $T$ before $T$ crosses the 0.10 crisis threshold. Scenario: step upgrade $\Delta F = 0.40$ from $F_0 = 0.80$ (above $\Delta F_{\mathrm{collapse}} = 0.31$), no workforce investment, starting from the surplus equilibrium. Lower $\beta$ (stickier customer trust) lengthens the warning window; $\sigma$ variation contributes less than two months at fixed $\beta$. Indirect-channel share: fraction of CLV uplift from $f$ via the sensing channel ($p{=}0$ counterfactual) at $F = 0.85$; values $<1\%$ confirm the capacity-multiplier channel dominates. See \S 6.1 of the main paper.}
\end{table}


\textit{CLV calibration.} The customer-lifetime-value mapping $\mathrm{CLV}(T) = w(T)/(1-\rho(T))$ uses illustrative parameters ($\rho_0 = 0.80$, $w_0 = 80$). Because both $\rho(\cdot)$ and $w(\cdot)$ are monotonically increasing in $T$, and $T^*$ itself is monotone in $f$ (Theorem~2), all qualitative predictions (sign of $\partial\,\mathrm{CLV}^*/\partial f$, existence of the Cliff, complementarity) are invariant to the CLV parametrisation. Only quantitative magnitudes (e.g., the 40\% CLV gap between scenarios) depend on these choices.

\section{Calibration Protocol}\label{sec:ec:calibration}

\subsection{Parameter Sources}

All time units are months. Rate parameters are expressed in months$^{-1}$; implied half-lives are $\ln(2)/\theta$. The calibration follows the partial-model testing and pattern-oriented fitting protocol of \citet{homer1993partial} and satisfies the validity criteria of \citet{barlas1996formal}. Three parameters warrant additional discussion beyond Table~\ref{tab:ec:paramsource}; the remaining seven are fully documented there.

\begin{table}
\TABLE
{Parameter calibration summary.\label{tab:ec:paramsource}}
{{\footnotesize\setlength{\tabcolsep}{3pt}%
\begin{tabular}{@{}lp{0.06\textwidth}cp{0.28\textwidth}p{0.26\textwidth}@{}}
\toprule
Parameter & Base & $\tau_{1/2}$ & Basis & Sensitivity \\
\midrule
$b$ (demand super-lin.) & 1.30 & -- & \citet{bell2018inventory}; \citet{acimovic2015fulfillment} & Trap requires $b>1$; robust \\
$a$ (trust--demand) & 0.40 & -- & \citet{mcknight2002developing}; \citet{sirdeshmukh2002consumer} & Weakens Trap at $a{=}0.2$; holds \\
$p$ (complementarity) & 0.50 & -- & Illustrative; \citet{krusell2000capital} & All 3 weaken at $p{=}0.25$ \\
$\sigma$ (overload asym.) & 2.00 & -- & \citeauthor{oliva2001cutting} (\citeyear{oliva2001cutting}) & Cliff disappears at $\sigma{=}1$ \\
$\eta$ (trust erosion asym.) & 2.00 & -- & \citet{sirdeshmukh2002consumer} & Weakens Trap at $\eta{=}1$ \\
$r$ (sensing-adapt.) & 0.40 & 1.7\,mo & \citet{tucker2007operational}; \citet{argote1990learning} & Robust \\
$\delta$ (depreciation) & 0.03 & 23\,mo & Illustrative & Robust ($\pm50\%$) \\
$\beta$ (trust-update) & 0.12 & 5.8\,mo & Illustrative & Primary lag driver (EC Tab.~EC.4) \\
$\mu$ (nat.\ $W$ decay) & 0.05 & 14\,mo & \citet{demerouti2001job}; \citet{dube2020minimum} & Robust \\
$\lambda$ (nat.\ $T$ decay) & 0.025 & 28\,mo & Illustrative & Robust \\
\bottomrule
\end{tabular}}}
{$\tau_{1/2}$: implied half-life $= \ln(2)/\theta$.}
\end{table}

\textbf{Demand super-linearity ($b = 1.3$).} A 10\% front-stage increase raises back-stage demand by approximately 13\%. The Trap requires only $b > 1$; sensitivity in Table~\ref{tab:ec:sensitivity}.

\textbf{Workforce-technology complementarity ($p = 0.5$).} A burned-out workforce ($W = 0$) delivers 50\% of design capacity \citep{krusell2000capital}. All results hold for any $p \in (0,1)$.

\textbf{Overload-recovery asymmetry ($\sigma = 2.0$).} Magnitude comparator from \citet{oliva2001cutting}; the Cliff disappears only at $\sigma = 1$ (Table~\ref{tab:ec:sensitivity}).

\subsection{Simulation Initial Conditions}

Table~\ref{tab:ec:simparams} reports initial conditions for each figure. All simulations use baseline parameters from Table~2 of the main paper, starting in the surplus regime ($\Keff > D$). Basin-of-attraction grid: $37{\times}26$ trajectories; doubled-grid validation ($73{\times}51$) shifts the separatrix by at most $0.0075$ $W_0$-units.

\begin{table}
\TABLE
{Simulation initial conditions and controls by figure.\label{tab:ec:simparams}}
{{\footnotesize\setlength{\tabcolsep}{2.5pt}%
\begin{tabular}{@{}lp{0.20\textwidth}ccccp{0.18\textwidth}@{}}
\toprule
Fig. & $(K_0, W_0, T_0)$ & $F_0$ & $u$ & $f$ & $t_{\max}$ & Notes \\
\midrule
2 & balance$^{\dagger}$ & varies & 0.05 & 0.03 & 500 & $37{\times}26$ grid; CLV thr.\ $=600$ \\
3 & \multicolumn{5}{c}{Analytical: $F_{\max}(W)=(K^* h(W))^{1/b}$; $K^*=u/\delta$} & Boundary formula \\
4 & $(1.60, 0.55, 0.50)$ & 0.80 & 0.05 & varies & 120 & Three scenarios (A,B,C) \\
5 & $(K^*\!,W^*\!,T^*)$ & 0.80 & --- & 0.01--0.10 & 500 & 61-pt $\Delta F$ sweep \\
EC & eq.~IC at $F{=}0.85$ & 0.85 & 0.05 & 0.00--0.20 & 120 & Study~3: Thm~2 decomp.; $p{=}0$ counterfactual confirms structural-zero \\
\bottomrule
\end{tabular}}}
{$^{\dagger}$Balance manifold: $K_0 h(W_0)=(1{+}a{\cdot}0.5)F^b$, $T_0{=}0.50$. Grid: $37{\times}26$ (962~pts). Collapse: $T(500){<}0.15$.}
\end{table}

\subsection{CLV Functional Form}

We model CLV using the geometric-series structure of \citet{gupta2004valuing}: $\mathrm{CLV} = w(T) / (1 - \rho(T))$, where $\rho(T) = 0.80 + 0.18T$ is an annual retention probability and $w(T) = 80(1 + 0.25\max(T-0.15, 0))$ an annualised per-customer margin. The retention endpoints ($\rho(0) = 0.80$, $\rho(1) = 0.98$) match the 70--90\% annual range in \citet{gupta2004valuing}; \$80 is an illustrative annual margin for US omnichannel grocery. The model's analytical results require only $\partial\,\mathrm{CLV}/\partial T^* > 0$, which holds everywhere since $\partial\,\mathrm{CLV}/\partial T^* \geq w(T^*)\rho'(T^*)/(1-\rho(T^*))^2 > 0$.


% ============================================================
% EC BIBLIOGRAPHY (subset of main paper references cited in EC)
% ============================================================

% -----------------------------------------------------------------------
% EC.5: Empirical Appendix
% -----------------------------------------------------------------------

\section{Empirical Appendix}\label{sec:ec:empirics}
{\small\setlength{\parskip}{2pt}
This appendix supplements the empirical analysis in \S5 of the main paper.

\subsection*{US Data and Robustness}

\textit{US data sources.}
ACSI scores (0--100, \texttt{theacsi.org}); Glassdoor ratings (1--5, min 5 reviews/firm-year); SEC EDGAR 10-K financials via XBRL API; FRED series ECOMPCTSA (e-commerce share); BLS JOLTS (turnover/wages).

\textit{Variables.}
$K{=}$CapEx/Revenue; $W{=}$Glassdoor overall; $T{=}$ACSI (0--100); $K{\times}W{=}$standardised interaction.

\textit{US merge.}
EDGAR backbone (276 obs, 17 firms, 2008--2024); left-join ACSI, Glassdoor, macro. Complete: 131 obs (47\%). No duplicate keys.

\textit{US robustness.}
(i)~Between-effects: $K{\times}W{=}+1.95$, $p{=}0.41$;
(ii)~Mundlak CRE;
(iii)~Lagged $K_{t-1}{\times}W_{t-1}$;
(iv)~COVID exclusion;
(v)~Leave-one-out;
(vi)~Wild bootstrap (9,999 reps);
(vii)~CR3 with $\mathrm{df}{=}G{-}k{=}9$.

\subsection*{Event-Study Pre-Trends}

Table~\ref{tab:ec:pretrends} reports event-time coefficients from the level DID and moderation specifications (\S5.2; $N = 108$, 16 firms). Pre-period coefficients ($t=-3$, $t=-2$) are insignificant in both specifications, supporting parallel trends.

\begin{table}[htbp]
\TABLE
{Event-time coefficients: parallel-trends check.\label{tab:ec:pretrends}}
{{\setlength{\tabcolsep}{4pt}\small
\begin{tabular}{@{}lrrrrrr@{}}
\toprule
& \multicolumn{3}{c}{\textit{Level DID}} & \multicolumn{3}{c}{\textit{Moderation ($\times W^z_{\mathrm{pre}}$)}} \\
\cmidrule(lr){2-4}\cmidrule(lr){5-7}
Event time & Coef. & $t$-stat & $p$-value & Coef. & $t$-stat & $p$-value \\
\midrule
$t = -3$ & $-0.125$ & $-0.40$ & 0.696 & $-2.093$ & $-0.11$ & 0.918 \\
$t = -2$ & $-0.375$ & $-1.15$ & 0.269 & $-2.174$ & $-0.11$ & 0.914 \\
$t = -1$ & \multicolumn{3}{c}{(reference)} & \multicolumn{3}{c}{(reference)} \\
$t = 0$ & $-0.367$ & $-2.21$ & 0.043 & $-1.266$ & $-0.06$ & 0.952 \\
$t = +1$ & $-0.167$ & $-0.65$ & 0.528 & $-0.876$ & $-0.04$ & 0.967 \\
$t = +2$ & $-0.300$ & $-0.78$ & 0.447 & $-0.875$ & $-0.04$ & 0.967 \\
$t = +3$ & $-0.900$ & $-1.69$ & 0.113 & $-1.033$ & $-0.05$ & 0.961 \\
\bottomrule
\end{tabular}
}}{Dependent variable: ACSI score. Firm FE; cluster-robust SEs with $p$-values from $t(15)$. Reference period: $t = -1$. Level DID: event-time dummies only. Moderation: interaction of event-time dummies with pre-shock standardised Glassdoor ($W^z_{\mathrm{pre}}$). All pre-period coefficients ($t = -3$, $t = -2$) are insignificant in both specifications.}
\end{table}

\subsection*{EU Replication (Exploratory)}

An independent panel covers 11 European retailers (UK, France, Germany, Spain, Sweden, Netherlands; 2019--2023; $N{=}20$ complete firm--year observations). Customer satisfaction uses Trustpilot (1--5); workforce well-being uses Glassdoor; financials from Yahoo Finance with within-year asset percentile rank as currency-neutral size control. E-commerce share from UK ONS (J4MC) and Eurostat (\texttt{isoc\_ec\_eseln2}).

\textit{EU results.}
The between-effects estimate of $\beta_3$ ($K{\times}W$) is positive ($+$0.47; CR2: $p = 0.045$ at $t(10)$). However, HC3 ($p = 0.53$) and Webb bootstrap \citep{webb2023reworking} ($p = 0.62$) do not reject zero, reflecting $N_c = 11$ clusters \citep{cameron2008bootstrap}. This sample is below standard thresholds for panel regression and is reported as an exploratory out-of-sample check.

\textit{EU robustness.}
(i)~Currency-neutral controls (assets rank);
(ii)~UK-only/Continental sub-samples;
(iii)~Leave-one-out: $K{\times}W \in [-0.04, +0.82]$;
(iv)~Wild bootstrap: $p{=}0.62$;
(v)~Imbens--Kolesar $\mathrm{df}_{\mathrm{eff}}{=}7.2$, $p{=}0.055$.
}% end \small

{\footnotesize\setlength{\bibsep}{1pt}
\begin{thebibliography}{99}

\bibitem[Acimovic and Graves(2015)]{acimovic2015fulfillment}
Acimovic J, Graves SC (2015) Making better fulfillment decisions on the fly in an online retail environment.
\textit{Manufacturing \& Service Operations Management} 17(1):34--51.

\bibitem[Argote and Epple(1990)]{argote1990learning}
Argote L, Epple D (1990) Learning curves in manufacturing.
\textit{Science} 247(4945):920--924.

\bibitem[Barlas(1996)]{barlas1996formal}
Barlas Y (1996) Formal aspects of model validity and validation in system dynamics.
\textit{System Dynamics Review} 12(3):183--210.

\bibitem[Bell et~al.(2018)]{bell2018inventory}
Bell DR, Gallino S, Moreno A (2018) Offline showrooms in omnichannel retail: Demand and operational benefits.
\textit{Management Science} 64(4):1629--1651.

\bibitem[Cameron et~al.(2008)]{cameron2008bootstrap}
Cameron AC, Gelbach JB, Miller DL (2008) Bootstrap-based improvements for inference with clustered errors.
\textit{Review of Economics and Statistics} 90(3):414--427.

\bibitem[Demerouti et~al.(2001)]{demerouti2001job}
Demerouti E, Bakker AB, Nachreiner F, Schaufeli WB (2001) The Job Demands-Resources model of burnout.
\textit{Journal of Applied Psychology} 86(3):499--512.

\bibitem[Dub{\'{e}} et~al.(2020)]{dube2020minimum}
Dub{\'{e}} A, Reich A, Shierholz H, West M (2020) Minimum wage shocks, employment flows and labor market frictions.
\textit{Journal of Labor Economics} 38(4):1009--1052.

\bibitem[Filippov(1988)]{filippov1988differential}
Filippov AF (1988) \textit{Differential Equations with Discontinuous Righthand Sides}.
Dordrecht: Kluwer Academic Publishers.

\bibitem[Gupta et~al.(2004)]{gupta2004valuing}
Gupta S, Lehmann DR, Stuart JA (2004) Valuing customers.
\textit{Journal of Marketing Research} 41(1):7--18.

\bibitem[Homer(1993)]{homer1993partial}
Homer JB (1993) Partial-model testing as a validation tool for system dynamics.
\textit{System Dynamics Review} 9(2):127--147.

\bibitem[Krusell et~al.(2000)]{krusell2000capital}
Krusell P, Ohanian LE, R\'ios-Rull JV, Violante GL (2000) Capital-skill complementarity and inequality: A macroeconomic analysis.
\textit{Econometrica} 68(5):1029--1053.

\bibitem[McKnight et~al.(2002)]{mcknight2002developing}
McKnight DH, Choudhury V, Kacmar C (2002) Developing and validating trust measures for e-commerce: An integrative typology.
\textit{Information Systems Research} 13(3):334--359.

\bibitem[Oliva and Sterman(2001)]{oliva2001cutting}
Oliva R, Sterman JD (2001) Cutting corners and working overtime: Quality erosion in the service industry.
\textit{Management Science} 47(7):894--914.

\bibitem[Sirdeshmukh et~al.(2002)]{sirdeshmukh2002consumer}
Sirdeshmukh D, Singh J, Sabol B (2002) Consumer trust, value, and loyalty in relational exchanges.
\textit{Journal of Marketing} 66(1):15--37.

\bibitem[Tucker(2007)]{tucker2007operational}
Tucker AL (2007) An empirical study of system improvement by frontline employees in hospital units.
\textit{Manufacturing \& Service Operations Management} 9(4):492--505.

\bibitem[Webb(2023)]{webb2023reworking}
Webb MD (2023) Reworking wild bootstrap-based inference for clustered errors.
\textit{Canadian Journal of Economics} 56(3):839--858.

\end{thebibliography}
}% end \small

\end{document}
